\begin{explanationbox}
	\textbf{\textcolor{CExplanation}{一句话直觉}}
	\par 平面内任意向量都可以被一组不共线的基底唯一线性表示，就像用两个独立方向作为坐标轴来描述位置。

	\medskip
	\textbf{\textcolor{CExplanation}{核心拆解}}
	\begin{description}[style=nextline, leftmargin=2.8em, labelsep=0.5em, itemsep=0.45em]
		\item[\textbf{要点一（基底不共线）：}]
			基底向量 $\vec{e}_1,\vec{e}_2$ 必须不共线，才能“张成”整个平面；若共线，则只能表示一条直线上的向量。
		\item[\textbf{要点二（系数唯一性）：}]
			给定不共线基底，系数 $\lambda_1,\lambda_2$ 由向量 $\vec{a}$ 完全确定，不同系数对应不同的向量。
	\end{description}

	\medskip
	\textbf{\textcolor{CExplanation}{几何本质（平行四边形法则）}}
	\par 以 $\vec{a}$ 为对角线，沿 $\vec{e}_1$ 和 $\vec{e}_2$ 方向作平行四边形，$\lambda_1\vec{e}_1$ 和 $\lambda_2\vec{e}_2$ 就是两个邻边。分解是否唯一的等价问题是：平行四边形是否唯一？由于 $\vec{e}_1,\vec{e}_2$ 不共线，邻边确定，从而分解唯一。

	\medskip
	\textbf{\textcolor{CExplanation}{代数意义（线性表示）}}
	\par $\vec{a} = \lambda_1\vec{e}_1 + \lambda_2\vec{e}_2$ 是向量的一种线性组合，它表示每个向量都是基底的线性函数。这种表示是建立向量坐标的理论基础。

	\medskip
	\textbf{\textcolor{CExplanation}{考点价值（参数求法与证共线）}}
	\begin{itemize}[leftmargin=*, itemsep=0.5em, topsep=0.4em]
		\item \textbf{考法一（求分解系数）：} 给定 $\vec{a}$ 与基底，通过方程或几何关系列式解出 $\lambda_1,\lambda_2$。
		\item \textbf{考法二（证三点共线）：} 若 $\vec{a}$ 能写成 $\mu\vec{b} + (1-\mu)\vec{c}$ 的形式，则终点共线，本质是基底分解的系数关系。
	\end{itemize}

	\medskip
	\textbf{\textcolor{CExplanation}{顿悟点}}
	\par\textbf{基底可以任意选取，但一旦选定，分解方式就是唯一确定的。不同基底下的系数不同，但向量本身不变。}

	\medskip
	\textbf{\textcolor{CExplanation}{使用场景}}
	\begin{itemize}[leftmargin=*, itemsep=0.5em, topsep=0.4em]
		\item \textbf{场景一（向量表达式运算）：} 涉及 $\overrightarrow{AB}$ 等难以直接计算的向量时，选取基底后用系数展开。
		\item \textbf{场景二（建系与坐标化）：} 优先选择垂直且模长已知的基底，使系数转化为坐标运算。
	\end{itemize}

\end{explanationbox}
